\section*{Capítulo \ref{ch:b1:p-block-13-15} --- O bloco p I: grupos do boro, do carbono e do nitrogênio}

\begin{solution}{exo:b1:p-block-13-15:1}
\ce{NH4NO3}: $-$III em \ce{NH4+}, $+$V em \ce{NO3-}. \ce{N2O5}: $+$V.
\ce{HNO2}: $+$III. \ce{N2H4}: $-$II. \ce{Mg3N2}: $-$III.
\end{solution}

\begin{solution}{exo:b1:p-block-13-15:2}
\ce{BF3}: boro com três ligações e seis elétrons, trigonal plano.
\ce{NH3}: três ligações e um \omterm{def:b1:lewis-resonance-vsepr:lewis}{par não ligante}. O \omterm{def:b1:lewis-resonance-vsepr:lewis}{par não ligante} de N forma
uma \omterm{def:b1:lewis-resonance-vsepr:lewis-acid}{ligação dativa} com B: \ce{BF3} é o \omterm{def:b1:lewis-resonance-vsepr:lewis-acid}{ácido de Lewis}, \ce{NH3} a \omterm{def:b1:lewis-resonance-vsepr:lewis-acid}{base de
Lewis}; no aduto, os dois átomos são tetraédricos.
\end{solution}

\begin{solution}{exo:b1:p-block-13-15:3}
\ce{CO2 + 2NaOH -> Na2CO3 + H2O}; \ce{SiO2 + 2NaOH -> Na2SiO3 + H2O};
\ce{P4O10 + 12NaOH -> 4Na3PO4 + 6H2O}; \ce{Al2O3 + 6HCl -> 2AlCl3 + 3H2O};
\ce{Al2O3 + 2NaOH + 3H2O -> 2NaAl(OH)4}.
\end{solution}

\begin{solution}{exo:b1:p-block-13-15:4}
$\Delta_r G^\circ = 2 \times (-16.45) = \qty{-32.9}{kJ/mol}$, $\log K =
32.9/5.708 = 5.76$, $K = 10^{5.76}$.
\end{solution}

\begin{solution}{exo:b1:p-block-13-15:5}
Anel de seis unidades \ce{SiO3^2-}: \ce{Si6O18^12-}. Cadeia dupla, para
quatro silícios: dois com dois vértices compartilhados (3 O, carga $-2$
cada), dois com três ($2.5$ O, carga $-1$ cada): \ce{Si4O11^6-}.
\end{solution}

\begin{solution}{exo:b1:p-block-13-15:6}
\qty{550}{kg} de \ce{Al(OH)3} são $550/78.0 = \qty{7.05}{kmol}$, que dão
\qty{3.53}{kmol} de \ce{Al2O3}, \qty{360}{kg}; com \qty{92}{\percent}:
\qty{331}{kg}.
\end{solution}

\begin{solution}{exo:b1:p-block-13-15:7}
O boro tem um orbital vazio e aceita o \omterm{def:b1:lewis-resonance-vsepr:lewis}{par não ligante} de um íon
hidróxido retirado da água: \ce{B(OH)3 + 2H2O <=> B(OH)4- + H3O+}. O
próton liberado vem de uma molécula de água, não do ácido bórico.
\end{solution}

\begin{solution}{exo:b1:p-block-13-15:8}
N passa de $-$III a $+$II, cinco elétrons por \ce{NH3}; cada \ce{O2} capta
quatro. $4 \times 5 = 5 \times 4$: \ce{4NH3 + 5O2 -> 4NO + 6H2O}.
\end{solution}

\begin{solution}{exo:b1:p-block-13-15:9}
\ce{NH3 + HNO3 -> NH4NO3}. Uma tonelada são $10^6/80.0 = \qty{1.25e4}{mol}$;
\qty{1.25e4}{mol} de amônia para neutralizar e \qty{1.25e4}{mol} para
produzir o ácido nítrico: $2.50 \times 10^4 \times 17.0 = \qty{425}{kg}$.
\end{solution}

\begin{solution}{exo:b1:p-block-13-15:10}
Os átomos do segundo \omterm{def:b1:periodicity:table}{período} são pequenos: seus orbitais \textit{p} se
sobrepõem bem lateralmente e formam ligações $\pi$ fortes, de modo que
\ce{N#N} e \ce{O=C=O} são moléculas pequenas e estáveis. Os átomos do
terceiro \omterm{def:b1:periodicity:table}{período} são maiores; sua sobreposição $\pi$ é fraca, e eles
preferem mais ligações simples: tetraedros \ce{P4}, a rede de \ce{SiO4}.
\end{solution}

\begin{solution}{exo:b1:p-block-13-15:11}
$E^\circ(\ce{PbO2}/\ce{Pb^2+}) = 1.46 > E^\circ(\ce{Cl2}/\ce{Cl-}) = 1.36$~V:
\ce{PbO2 + 4H+ + 2Cl- -> Pb^2+ + Cl2 + 2H2O}. O chumbo(IV) é instável em
relação ao chumbo(II) (par inerte); o silício(II) não é um estado estável,
e \ce{SiO2} não é \omterm{def:b1:nernst:couple}{oxidante}.
\end{solution}

\begin{solution}{exo:b1:p-block-13-15:12}
$160 \times 17.0/14.0 = \qty{194}{Mt}$ de amônia, que contêm $194 \times
3.0/17.0 = \qty{34}{Mt}$ de hidrogênio, produzido principalmente a partir
do gás natural (metano e vapor d'água).
\end{solution}

\begin{solution}{pb:b1:p-block-13-15:1}
\textbf{1.} :N\,$\equiv$\,N: --- uma ligação tripla, muito forte, sem polaridade.
\textbf{2.} 0, $-$III, $+$II, $+$IV, $+$V, e $-$III/$+$V em \ce{NH4NO3}.
\textbf{3.} Bactérias nas raízes das leguminosas, os raios, o processo
Haber--Bosch.
\textbf{4.} Romper a ligação tripla exige um \omterm{def:b1:elementary-steps:catalyst}{catalisador} que as plantas
não têm; as bactérias com a enzima nitrogenase têm.
\textbf{5.} O nitrogênio é reduzido (de 0 a $-$III), o hidrogênio oxidado (de 0 a $+$I).
\textbf{6.} \ce{N2 + 3H2 <=> 2NH3}.
\textbf{7.} $K = 10^{2 \times 16.45/5.708} = 10^{5.76}$.
\textbf{8.} A \qty{25}{\celsius} a reação é lenta demais; o \omterm{def:b1:elementary-steps:catalyst}{catalisador} só
funciona a quente, e as condições são um compromisso discutido no volume
do 2.º~ano.
\textbf{9.} \qty{58.8}{kmol} de amônia: \qty{29.41}{kmol} de \ce{N2},
\qty{824}{kg}; \qty{88.2}{kmol} de \ce{H2}, \qty{176}{kg}.
\textbf{10.} $V(\ce{N2}) = 29.4 \times 10^3 \times 8.314 \times 298.15/10^5 =
\qty{729}{m^3}$; ar $729/0.78 = \qty{935}{m^3}$.
\textbf{11.} Do gás natural (reforma a vapor do metano), com \ce{CO2} como
subproduto.
\textbf{12.} Injetada no solo, é um fertilizante barato, mas é um gás
tóxico; os sólidos (ureia, nitrato de amônio) são mais fáceis de
armazenar e espalhar.
\textbf{13.} N: $-$III $\to$ $+$II (5 elétrons), O: 0 $\to$ $-$II (4 por
\ce{O2}): \ce{4NH3 + 5O2 -> 4NO + 6H2O}.
\textbf{14.} \ce{2NO + O2 -> 2NO2}.
\textbf{15.} \ce{3NO2 + H2O -> 2HNO3 + NO}: um \omterm{def:b1:e-ph-diagrams:disproportionation}{desproporcionamento} (N de
$+$IV a $+$V e $+$II).
\textbf{16.} \ce{NH3 + 2O2 -> HNO3 + H2O}.
\textbf{17.} $2 \times 58.8 \times 32.0 = \qty{3.76}{t}$.
\textbf{18.} Ela torna a oxidação rápida e seletiva para NO, em vez do
\ce{N2}, mais estável.
\textbf{19.} \ce{NH3 + HNO3 -> NH4NO3}.
\textbf{20.} Metade.
\textbf{21.} \qty{29.4}{kmol} de \ce{NH4NO3}, \qty{2.35}{t}.
\textbf{22.} $28.0/80.0 = \qty{35}{\percent}$.
\textbf{23.} Ele contém um \omterm{def:b1:nernst:couple}{oxidante} (o nitrato) e um \omterm{def:b1:nernst:couple}{redutor} (o amônio) no
mesmo \omterm{def:b1:crystals-metals:crystal}{cristal}: aquecido ou misturado a combustíveis, pode se decompor de
modo explosivo.
\textbf{24.} $58.8 \times 63.0 =$ \textbf{\qty{3.7}{t} de ácido nítrico por
tonelada de amônia}.
\end{solution}
